\scanpage{367}
\mspage{106}

comme contrafoncteur en $S$,
\[
H:S\longmapsto\pi_0\,\underline{\mathscr H\mathrm{om}}_T(S_T,\unsure{\mathcal C|C})\qquad(\mathrm{Cat})^\circ\longrightarrow(\mathrm{Ens}),
\]
et à une éventuelle expression de ce \unsure{cfoncteur|foncteur} $H$ en termes du ind-objet $H_C$ de $\mathrm{Cat\,hot}$, par la formule hypothétique
\[
H(S)\simeq\mathrm{Hom}_{\inserted{\mathrm{Ind}(\mathrm{Cat\,hot})}}(\gamma(S),H_C),
\]
où
\[
\gamma:\mathrm{Cat}\longrightarrow\mathrm{Cat\,hot}
\]
est le foncteur canonique. Si on avait un tel isomorphisme \inserted{canonique} fonctoriel en $S$, cela donnerait une description (à isom unique près) du ind-objet $H_C$ en termes du foncteur précédent $H$ sur $\mathrm{Cat}$. \deleted{\unsure{I}}

Il me vient un doute ici, en vue de l'exemple $\deleted{\unsure{C}}T=\underline{\mathrm{Ty}}$, $C=\underline{\mathrm{Ch}}_\Sigma(X;a,b)$. Les trois types d'objets a) b) c) ont le $\unsure{=|\text{même}}$ $\pi_0$, et il est sûr que ``c'est le bon''. D'autre part, pour une catégorie \unsure{paramètre|paramétrante} variable $S$, il est sûr aussi que le \unsure{bon} foncteur en $S$ est
\[
S\longmapsto\pi_0\bigl(\underline{\mathrm{Hom}}_T(S_T,\mathcal C)\bigr)\simeq\mathrm{Hom}_{\mathrm{Ind\,Cat\,hot}}(S,H_C)
\]
isomorphe aussi, \uline{si $\deleted{\unsure{D|S}}\unsure{S|D}$ est finie}, à
\[
\deleted{\unsure{\pi_0}}\mathrm{Hom}_{\mathrm{Cat\,hot}}(S,C_\infty),
\]
ce qui montre que, du moins pour des arguments $S\in(\mathrm{Cat})$ qui sont des cat. \uline{finies}, les objets

% archive sheet number: 366
